\subsubsection{Creación de conjunto de datos médicos de alcance intermedio (1 año)}

Un proceso similar al de la creación del árbol familiar fue aplicado a construir una historia clínica de Norman con especial enfoque en la correcta estructuración de las fechas (formato ISO 8601), el tipo de dolencia, la medicación prescrita, hospital visitado y cualquier detalle de relevancia para asemejar un registro real.

Se continúa con la normalización al idioma inglés como decisión de diseño, orientada a aprovechar que los modelos empleados fueron entrenados mayoritariamente sobre texto en inglés; en la práctica, esto tiende a reducir el consumo de tokens y a favorecer el rendimiento del procesamiento.

El resultado, que abarca el período de abril de 2023 a marzo de 2024, se presenta a continuación.

\paragraph{Enero 2024 (01-2024.json)}

El listado completo figura en el Anexo técnico (p.~\pageref{anexo:code-4-3-2-1}).

\textit{Nota:} Código disponible en \url{https://github.com/aditzend/hyperpersonalization/blob/main/app/notebooks/28-medical-graph.ipynb}

\paragraph{Febrero 2024 (02-2024.json)}

El listado completo figura en el Anexo técnico (p.~\pageref{anexo:code-4-3-2-2}).

\paragraph{Marzo 2024 (03-2024.json)}

El listado completo figura en el Anexo técnico (p.~\pageref{anexo:code-4-3-2-3}).

\paragraph{Abril 2023 (04-2023.json)}

El listado completo figura en el Anexo técnico (p.~\pageref{anexo:code-4-3-2-4}).

\paragraph{Registros médicos adicionales}

Los meses restantes (mayo a diciembre de 2023, con excepción de septiembre, que no registró eventos médicos) siguen un patrón similar de estructuración de datos, incluyendo:

\begin{itemize}
\item Mayo 2023: Consulta endocrinológica para función tiroidea y terapia física rutinaria
\item Junio 2023: Seguimiento para manejo de colesterol y mamografía de detección
\item Julio 2023: Evaluación neurológica por episodios de mareo y colonoscopía rutinaria
\item Agosto 2023: Vacunación anual contra la gripe y consulta ortopédica por dolor de rodilla
\item Octubre 2023: Consulta por dolores de cabeza persistentes y análisis de laboratorio rutinarios
\item Noviembre 2023: Examen ocular rutinario y sesión de terapia física
\item Diciembre 2023: Consulta por alergias estacionales y vacuna anual contra la gripe
\end{itemize}

Cada registro médico mantiene la misma estructura de campos: fecha y hora en formato ISO 8601, descripción, tipo de visita, especialidad médica, nombre del doctor, síntomas experimentados, hospital o centro de atención, medicamentos prescritos, terapias recibidas y notas adicionales.

\textit{Nota:} El conjunto completo de datos médicos está disponible en \url{https://github.com/aditzend/hyperpersonalization/blob/main/app/notebooks/28-medical-graph.ipynb} 